\pagebreak

\section{Place Groups}
\label{sec: place_groups}
\index{affinity!place groups}
\index{environment variables!OMP_PLACES@\kcode{OMP_PLACES}}
\index{OMP_PLACES@\kcode{OMP_PLACES}}

The \kcode{OMP_PLACES} environment variable is used to define the initial place group
and its associated place partition. Each place in the place group is an ordered
set of indivisible execution units (typically hardware threads) and confines
where a bound OpenMP thread is executed.

The specification provides a grammar that defines the format of strings to which
the environment variable can be set. The string can either be a numeric
place-group specification, formed from non-negative numbers, or an abstracted
place-group specification, formed from abstract names.

A numeric place-group specification is built up from non-negative numbers that
each represent the smallest unit of execution that the implementation supports,
such as a hardware thread. Which specific execution unit a given number
represents is implementation defined. Here are some examples where the numeric
specification is a simple list of these numbers:

\begin{boxeducode}
# 4 places corresponding, in order, to execution units:
# 0,1,2,3
\kcode{OMP_PLACES}="0,1,2,3"

# 4 places corresponding, in order, to execution units:
# 2,1,0,3
\kcode{OMP_PLACES}="2,1,0,3"

# 4 places corresponding, in order, to execution units:
# 0,2,4,6
\kcode{OMP_PLACES}="0,2,4,6"
\end{boxeducode}

An interval notation
\plc{lower-bound}\code{:}\plc{length}[\code{:}\plc{stride}] may be used
express an ordered sequence of numbers:

\begin{boxeducode}
# 4 places corresponding, in order, to execution units:
# 0,1,2,3
\kcode{OMP_PLACES}="0:4"

# 4 places corresponding, in order, to execution units:
# 0,2,4,6
\kcode{OMP_PLACES}="0:4:2"
\end{boxeducode}

An exclusion operator \code{!} may also be used to exclude a given number from a
sequence otherwise specified. If no sequence is otherwise specified, then the
exclusion operator is interpreted as expanding to all numbers that could be
specified except for the one that is excluded.

\begin{boxeducode}
# 7 places corresponding, in order, to execution units:
# 0,1,2,3,4,5,7 (6 excluded)
\kcode{OMP_PLACES}="0:8,!6"

# Assuming 16 total execution units, 15 places, corresponding,
# in order, to execution units:
# 0,1,2,3,4,5,7,8,9,10,11,12,13,14,15 (6 excluded)
\kcode{OMP_PLACES}="!6"
\end{boxeducode}

In the above examples, each place in the specified place group corresponded to a
single execution unit denoted by a non-negative number. A numeric specification
may group a subsequence of non-negative numbers in curly braces to indicate that
the subsequence collectively forms a place within the place group:

\begin{boxeducode}
# 4 places corresponding, in order, to ordered sets of
# execution units: {0,1},{2,3},{4,5},{6,7}
\kcode{OMP_PLACES}="{0,1},{2,3},{4,5},{6,7}"

# 4 places corresponding, in order, to ordered sets of
# execution units: {0},{1},{2,3,4},{5,6,7}
\kcode{OMP_PLACES}="0,1,{2,3,4},{5,6,7}"
\end{boxeducode}

The interval notation and exclusion operator may also be used as shorthand to
specify a subsequence of execution units that comprise a place:

\begin{boxeducode}
# 2 places corresponding, in order, to ordered sets
# of execution units: {0,1,2,3},{4,5,6,7}
\kcode{OMP_PLACES}="{0:4},{4:4}"

# 2 places corresponding, in order, to ordereded sets
# execution units: {0,2,4,6},{1,3,5,7}
\kcode{OMP_PLACES}="{0:4:2},{1:4:2}"

# 2 places corresponding, in order, to ordered sets
# execution units: {0,2,6},{8,12,14}
\kcode{OMP_PLACES}="{0:4:2,!4},{8:4:2,!10}"
\end{boxeducode}

In the above examples, the place groups were specified as a flat list of places
(which each consist of one or more execution units corresponding to non-negative
numbers). A place group may also consist of nested place groups, such that a
place group can be visualized as a tree (nodes with children being place groups
and leaf nodes being places). Angle brackets are used to specify nested place
groups within a place group specification:

\begin{boxeducode}
# 8 places corresponding, in order, to execution units:
# 0,1,2,3,4,5,6,7
# place-group-var = group 1, group 2
#         group 1 = 0,1,2,3
#         group 2 = 4,5,6,7
\kcode{OMP_PLACES}="<0,1,2,3>,<4,5,6,7>"

# 5 places corresponding, in order, to ordered sets of
# execution units: {0},{1,2},{4},{5},{6,7}
# place-group-var = 0, {1,2}, group 1
#         group 1 = 4,5,{6,7}
\kcode{OMP_PLACES}="0,{1,2},<4,5,{6,7}>"

# 7 places corresponding, in order, to ordered sets of
# execution units: {0},{1},{2},{3},{4,5},{6},{7}
# place-group-var = group 1, {4,5}, group 2
#         group 1 = group 3, group 4
#         group 2 = 6,7
#         group 3 = 0,1
#         group 4 = 2,3
\kcode{OMP_PLACES}="<<0,1>,<2,3>>,{4,5},<6,7>"
\end{boxeducode}

The interval notation and exclusive operator may also be used to specify a list
of place groups:

\begin{boxeducode}
# 16 places corresponding, in order, to execution units 
# 0, 1, 2, 3, 4, 5, 6, 7, 8, 9,
# 10, 11, 12, 13, 14, and 15
# place-group-var = group 1, group 2, group 3, group 4
#         group 1 = 0,1,2,3
#         group 2 = 4,5,6,7
#         group 3 = 8,9,10,11
#         group 4 = 12,13,14,15
\kcode{OMP_PLACES}="<0:4>:4:4"

# 16 places corresponding, in order, to execution units 
# 0, 1, 2, 3, 4, 5, 6, 7, 12, 13,
# 14, 15, 16, 17, 18, and 19
# place-group-var = group 1, group 2, group 3, group 4
#         group 1 = 0,1,2,3
#         group 2 = 4,5,6,7
#         group 3 = 12,13,14,15
#         group 4 = 16,17,18,19
\kcode{OMP_PLACES}="<0:4>:2:4,<12:4>:2:4"

# 16 places corresponding, in order, to execution units 
# 0, 1, 2, 3, 4, 5, 6, 7, 12, 13,
# 14, 15, 16, 17, 18, and 19
# place-group-var = group 1, group 2, group 3, group 4
#         group 1 = 0,1,2,3
#         group 2 = 4,5,6,7
#         group 3 = 12,13,14,15
#         group 4 = 16,17,18,19
\kcode{OMP_PLACES}="<0:4>:5:4,!<8:4>"
\end{boxeducode}

An abstracted place-group specification uses place-group abstract names to
specify a place group. Which specific execution units these abstract names
represent is implementation defined. Predefined abstract names are available
that have particular meanings (as interpreted by the implementation):

\begin{compactitem}
\item \kcode{threads}: each place is a hardware thread
\item \kcode{cores}: each place is a core
\item \kcode{ll_caches}: each place is a set of cores for a single last-level cache
\item \kcode{numa_domains}: each place is a set of cores of a single NUMA domain
\item \kcode{sockets}: each place is a set of cores for a single socket
\end{compactitem}

Here are some examples where a simple place-group abstract name is used to
specify the place group:

\begin{boxeducode}
# Assuming the device supports 16 hardware threads,
# 16 places corresponding, in some order, to execution
# units that correspond to hardware threads:
# 0,8,1,9,2,10,3,11,4,12,5,13,6,14,7,15
\kcode{OMP_PLACES}="threads"

# Assuming the device supports 8 cores of 2 hardware threads,
# 8 places corresponding, in some order, to ordered sets
# of execution units that correspond to cores:
# {0,8},{1,9},{2,10},{3,11},{4,12},{5,13},{6,14},{7,15}
\kcode{OMP_PLACES}="cores"

# Assuming the device supports 2 sockets of 8 hardware threads,
# 2 places corresponding, in some order, to ordered sets
# of execution units that correspond to sockets:
# {0,8,1,9,2,10,3,11},{4,12,5,13,6,14,7,15}
\kcode{OMP_PLACES}="sockets"
\end{boxeducode}

An abstract name can be augmented with a length and optional stride enclosed in
parentheses with the syntax \code{(}\plc{len}[\code{:}\plc{stride}]\code{)}
to specify that only a subset of the execution units should be selected for
inclusion in the place group:

\begin{boxeducode}
# Assuming the device supports at least 8 hardware threads,
# the first 8 places of a list (0,8,1,9,...) corresponding,
# in some order, to execution units that correspond to hardware
# threads: 0,8,1,9,2,10,3,11
\kcode{OMP_PLACES}="threads(8)"

# Assuming the device supports at least 16 hardware threads,
# the first 8 places with a stride of 2 of a list (0,8,1,9,...)
# corresponding, in some order, to execution units that
# correspond to hardware threads: 0,1,2,3,4,5,6,7
\kcode{OMP_PLACES}="threads(8:2)"

# Assuming the device supports at least 4 cores,
# the first 4 places of a list ({0,8},{1,9},...) corresponding,
# in some order, to ordered sets of execution units that
# correspond to cores:
# {0,8},{1,9},{2,10},{3,11}
\kcode{OMP_PLACES}="cores(4)"

# Assuming the device supports at least 8 cores,
# the first 4 places with a stride of 2 of a list
# ({0,8},{1,9},...) corresponding, in some order, to ordered
# sets of execution units that correspond to cores:
# {0,8},{2,10},{4,12},{6,14}
\kcode{OMP_PLACES}="cores(4:2)"
\end{boxeducode}

The abstract name may alternatively be appended with an index list enclosed in
square brackets. Each item of the index list is an index, an index interval, or
an excluded index. An index refers to a place within the place group given by
the abstract name.

\begin{boxeducode}
# 3 places corresponding, in order, to execution units
# that correspond to some hardware threads: 0,2,3
\kcode{OMP_PLACES}="threads[0,2,3]"

# 7 places corresponding, in order, to execution units
# that correspond to some hardware threads:
# 0,2,4,6,1,5,7 (3 excluded)
\kcode{OMP_PLACES}="threads[0:4:2,1:4:2,!3]"
\end{boxeducode}

An abstract name may also specify a nested abstract name for more control over
execution unit selection. Of the predefined place-group abstract names, threads
may be nested within cores, and both cores and threads may be nested within the
other abstract names. No abstract name may be nested inside threads, and only
threads may be nested inside cores. Whether any other nestings are allowed is
considered implementation defined.

Two types of nesting is possible with abstract names. The dot operator
\ucode{.} treats the combined abstract name as referring to the sets of
execution units given by the nested abstract name in the context of the outer
abstract name. Examples:

\begin{boxeducode}
# Assuming the device supports 8 cores and at least 2 hardware
# threads per core, 16 places corresponding to 2 hardware threads
# within each of the 8 cores:
# 0,8,1,9,2,10,3,11,4,12,5,13,6,14,7,15
\kcode{OMP_PLACES}="cores.threads(2)"

# Assuming the device supports 8 cores and at least 2 hardware
# threads per core, 4 places corresponding to 2 hardware
# threads within each of 2 cores:
# 0,8,1,9
\kcode{OMP_PLACES}="cores(2).threads(2)"

# Assuming the device supports 8 cores and 4 hardware
# threads per core, 8 places corresponding to the first 2
# hardware threads on cores 1, 3, 5, and 7:
# 1,9,3,11,5,13,7,15
\kcode{OMP_PLACES}="cores[1:4:2].threads[0:2]"

# Assuming the device supports 2 sockets, each with 8 cores
# and 2 hardware threads per core, 16 places corresponding to
# all hardware threads executing on cores of first socket:
# 0,16,1,17,2,18,3,19,4,20,5,21,6,22,7,23,8,24
\kcode{OMP_PLACES}="sockets[0].threads"

# Assuming the device supports 2 sockets, each with 8 cores
# and 2 hardware threads per core, 4 places corresponding,
# in some order, to 4 cores of second socket:
# {9,25},{10,26},{11,27},{12,28}
\kcode{OMP_PLACES}="sockets[1].cores(4)"

# Assuming the device supports 2 sockets, each with 8 cores
# and 2 hardware threads per core, 8 places corresponding
# to hardware threads of 4 cores of second socket:
# 9,25,10,26,11,27,12,28
\kcode{OMP_PLACES}="sockets[1].cores(4).threads"
\end{boxeducode}

An abstract name can also be specified with a nested abstract name enclosed in
angle brackets. This indicates that each set of execution units given by the
outer abstract name becomes a specified place group consisting of the sets of
execution units given by the nested abstract name.

\begin{boxeducode}
# Assuming the device supports 8 cores each with 2
# hardware threads, 16 places corresponding, in order,
# to execution units:
# 0,8,1,9,2,10,3,11,4,12,5,13,6,14,7,15
# place-group-var = group 1, group 2, group 3, group 4,
#                   group 5, group 6, group 7, group 8
#         group 1 = 0,8
#         group 2 = 1,9
#         group 3 = 2,10
#         group 4 = 3,11
#         group 5 = 4,12
#         group 6 = 5,13
#         group 7 = 6,14
#         group 8 = 7,15
\kcode{OMP_PLACES}="cores<threads>"

# Assuming the device supports 8 cores each with 2
# hardware threads, 4 places corresponding to 
# 2 hardware threads from 4 cores:
# 0,8,1,9,2,10,3,11
# place-group-var = group 1, group 2, group 3, group 4
#         group 1 = 0,8
#         group 2 = 1,9
#         group 3 = 2,10
#         group 4 = 3,11
\kcode{OMP_PLACES}="cores(4)<threads(2)>"
\end{boxeducode}
